\documentclass{article}

\title{Runge-Kutta Methods and Mechanical Oscillators}
\author{Jaden Krause}
\date{October 27, 2026}
\usepackage{float}
\usepackage{graphicx}
\usepackage{amsmath}
\makeatletter
\renewcommand{\topfraction}{0.95}       % Allow figures to use up to 95% of the page top
\renewcommand{\textfraction}{0.05}      % Allow text to take as little as 5% of the page 
\renewcommand{\bottomfraction}{0.95}    % Allow figures to use up to 95% of the page bottom
\renewcommand{\floatpagefraction}{0.85}  % Allow a float page to be created even if it's only 20% full
\setcounter{topnumber}{4}               % Allow up to 4 figures at the top of a page
\setcounter{bottomnumber}{4}            % Allow up to 4 figures at the bottom of a page
\setcounter{totalnumber}{8}             % Allow up to 8 total figures per page
\makeatother
\begin{document}
\maketitle

\section{Introduction}
Continuing the study of differential equations, it behooves us to find a numerical approximation method 
which yields more accuracy than Euler's Method. As seen in previous work, Euler's Method has inherent error 
due to its approach of finding the slope of the line at the starting point of the interval. There are 
ways we can compensate Euler's Method to yield greater accuracy, and there are other methods with higher 
orders of accuracy that are better suited to many problems. In this paper, we will consider these ramifications. 
\par One practical application of higher-order approximation methods is mechanical oscillators. This is 
also an object of interest when studying differential equations, as the functions for oscillators are 
derived using differential equations. More accurate approximation methods are needed for these problems 
because the slope changes frequently, and often faster than a method like Euler's could compensate for. 
In this paper, we will be examining how new methods of numerical approximation can be applied to oscillator 
problems.

\section{Better Methods of Approximation}
\subsection{Improving Euler's Method}
When we evaluated Euler's method, we found that it approximated the function by estimating the slope of the 
function at a point, moving forward a step size $h$, and repeating the process. This resulted in the approximation 
deviating from the true function, generally erring behind the curve and changing slower, ultimately shown 
in Euler's order of convergence being 1 ($\mathcal{O}(h)$). This tendency is magnified 
with graphs that change rapidly and large step size $h$. For a more in-depth discussion of these tendencies, refer 
to Krause, Jaden \textit{Initial Value Problems and Euler's Method}, 2026. 
\par To improve Euler's method and minimize these error tendencies, we employ what is referred to as the midpoint 
method, or Runge-Kutta 2 (RK2). This method takes the form: 
\begin{align*}
x_{n+1} &= x_n + h, \\
y_{n+1} &= y_n + hf(x_n + \frac{h}{2}, y_n + \frac{h}{2}f(x_n, y_n)).
\end{align*}
This method has convergence of order 2 ($\mathcal{O}(h^2)$). Practically, this method uses a similar approach to 
Euler's but approximates the slope of the function at the midpoint between the present time step $x_n$ and the 
next step at $x_{n+1}$ to obtain greater accuracy and follow the graph closer. 
\subsection{Higher-Order Convergence Methods}
With a little extra work, the midpoint formula can be modified to increase the convergence of the formula. This is 
done using the Taylor expansion; 
\begin{align*}
x_{n+1} &= x_n + h \\
y_{n+1} &= y_n + hf(x_n, y_n) + \frac{h^2}{2!}f_2(x_n, y_n) + \cdots + \frac{h^p}{p!}f_p(x_n, y_n)
\end{align*}
with convergence of order $\mathcal{O}(h^p)$ for a Taylor method of order $p$. This leaves us with the question, how 
many terms do we include? Accuracy increases with each further step, but computational difficulty increases exponentially 
as we need to compute more partial derivatives to obtain $f_p$ (full explanation of process and proofs found in Nagle, et al. 
132-133). Carrying out this formula to the fourth degree yields a convergence rate of $\mathcal{O}(h^4)$ with a formula that 
is easy to program and carry out. This is the most popular of the numerical approximation methods, and is known as Runge-Kutta 
4 (RK4):
%RUNGE-KUTTA 4 METHOD FORMULAS. ENSURE IN THE RIGHT PLACE OF DOCUMENT
\begin{align*}
    x_{n+1} &= x_n + h, \\
    y_{n+1} &= y_n + \frac{1}{6}(k_1 + 2k_2 + 2k_3 +k_4), 
\end{align*}
where
\begin{align*}
    k_1 &= hf(x_n, y_n), \\
    k_2 &= hf(x_n + \frac{h}{2}, y_n + \frac{k_1}{2}), \\
    k_3 &= hf(x_n + \frac{h}{2}, y_n + \frac{k_2}{2}), \\
    k_4 &= hf(x_n + h, y_n + k_3)
\end{align*}
Runge-Kutta 4's balance of accuracy and computational ease are what make it such a widespread solution for numerical integration 
in differential equations. Figure 1 shows an accuracy comparison between RK2 and RK4, in which RK4 is seen to adhere much closer 
to the true curve of the function. We will be using RK4 for the rest of this discussion as we turn to evaluating 
mechanical oscillators.
\begin{figure}
    \centering
    \includegraphics[width=1\textwidth]{FiguresofRK4/RK2vsRK4.png}
    \caption{Comparison between RK2 and RK4}
\end{figure}
\section{Mechanical Oscillators and Damping}
\subsection{The Problem}
In our study of differential equations, we are trying to build skills to evaluate the one-dimensional elasticity of metals. 
Having been introduced to RK4, it is now within our grasp to accurately evaluate mechanical oscillator systems such as springs. 
The equations for this involve unknown functions of $t$, $y''$ and $y'$, which is where RK4 applies itself. The particular problem we 
will use in this analysis is:
$$y'' + by' + 16y = 0, \qquad y(0) = 1, \qquad y'(0) = 0$$
for
$$ b = 0, \quad 6, \quad 8, \quad 10.$$
This equation follows the form:
$$my'' + by' + ky = 0$$
$m$ represents the mass on the end of the spring (measured in kg). $b$ is the damping coefficient, representing forces such as 
friction and air resistance (measured in kg/s or N*s/m). $k$ represents the spring constant (measured in N/m). The differential 
equation is equal to 0 because there is no external input of force, such as external vibration, in our example. The initial value 
condition for $y$ is equal to 1 at $t = 0$, showing that the mass starts 1 meter from the spring's equilibrium. The initial 
value condition for $y'$ is equal to 0 at $t = 0$, showing that the mass starts with no initial velocity. When evaluated at the 
indicated values of $b$ using RK4, the results in Figure 2 are obtained. 
\begin{figure}
    \centering
    \includegraphics[width=0.90\textwidth]{FiguresofRK4/VariousDampers.png}
    \caption{Damping Performance at Several $b$ Values}
\end{figure}
\subsection{Results}
When $b = 0$, the system is undamped. With no dissipation of force, the spring continues oscillating with no reduction in amplitude. 
When $b = 6$, the system is underdamped. This is because, while the spring is returned to equilibrium, it oscillates past equilibrium 
and has to spring back the other way to reach rest. This kind of uncontrolled bobbing is exemplified in weak vehicle suspension, where 
a irregularity in the road is hit and the vehicle continues moving up and down for a long time. When $b = 8$, the system is critically 
damped, where the system returns to equilibrium quickly without overcorrecting. When $b = 10$, the system is overdamped, experiencing 
so much friction that it doesn't return to equilibrium for a long time. An oscillator's damping behavior is determined by the 
discriminant $b^2 - 4mk$. If the expression is less than 0, the system is underdamped. If the expression is equal to 0, the system is 
critically damped. If the expression is greater than 0, the system is overdamped. 
\subsection{Conversion to a First-Order System}
In our example, the differential equation includes both the first and second derivatives of $y(t)$. This poses a problem for RK4, 
which can only estimate the change for a first-degree derivation of the function. This requires that we convert our system to first-
order. Mathematically, this is done by assigning new variables to our functions. Recall our original ODE:
$$y'' + by' + 16y = 0$$
This equation, being 2\textsuperscript{nd} order, requires 2 variables. Set $x_1$ equal to position, yielding $x_1 = y$. Then, set 
$x_2$ equal to velocity. This leads to $x_2 = x_1' = y'$. Thus, our first equation in the system is $x_2 = x_1'$. If $x_2 = y'$, it 
follows that $x_2' = y''$. If we rewrite our original equation to isolate $y''$, we get:
$$y'' = -16y - by'$$
Substituting in the variables we defined earlier yields:
$$x_2' = -16x_1 - bx_1', \qquad x_2 = x_1'$$
This equation is now in a form that can be dealt with using RK4.

\section{Convergence of Undamped Systems}
All this work that we've performed using RK4 begs the question: is this helpful or necessary? Would Euler's Method suffice to evaluate 
a mechanical oscillator? After all, it is simpler to use. The answer to these questions is that RK4 is absolutely necessary. Euler's Method 
overestimates a function because it estimates the slope at point $x$ and then moves a distance $h$ in a straight line. This tendency 
accentuates error when slope changes rapidly or when it changes from positive to negative. In an undamped oscillatory system, the slope 
changes constantly in this way. Ultimately, using Euler's Method to evaluate such a system results in accumulating exaggerations of the spring's 
position, implying that there is energy input to the system. We specified earlier that this is not the case. Moreover, the errors of Euler's 
Method continue to accumulate exponentially as time ($t$) approaches infinity. This explosion of error is unacceptable for a robust analysis 
of oscillatory systems. RK4, however, doesn't suffer from this tendency. In fact, it follows the function quite closely. This is reflective of 
how RK4 approximates the slope of the function at 4 points along the distance $h$. Figure 3 demonstrates both the accuracy of RK4 and the 
error accumulation of Euler's Method. Figure 4 is a log-log plot of the error for each method at various step sizes $h$, and demonstrates 
the first-order convergence of Euler's Method and the fourth-order convergence of RK4.
\begin{figure}[H]
    \centering
    \includegraphics[width=0.90\textwidth]{FiguresofRK4/EulervsRK4.png}
    \caption{Accuracy of Approximation Methods}
\end{figure}
\begin{figure}[H]
    \centering
    \includegraphics[width=0.90\textwidth]{FiguresofRK4/loglogEuVSRK4.png}
    \caption{Convergence of Approximation Methods}
\end{figure}

\newpage
\section{Conclusion}
In this paper, we discussed numerical integration methods with greater accuracy that Euler's Method. Specifically, we looked at the midpoint 
method, or RK2, and Runge-Kutta 4, or RK4. RK4 is the most popular of the numerical integration methods, due to its accuracy and relative 
computational ease. This method applies itself especially well in the example of mechanical oscillators, an application which Euler's Method 
fails to succeed at. Being able to approximate mechanical oscillators allows us to look at the different damping cases and to analyze the overall 
behavior of these systems. 

\section{Bibliography}
\noindent\hangindent=1.5em \hangafter=1
Nagle, R. Kent, Edward B. Saff, and Arthur David Snider (2018). \textit{Fundamentals of Differential Equations}. 9th ed. Pearson. \\
\hspace*{1.5em}\textsc{isbn}: 9780321977113

\noindent\hangindent=1.5em \hangafter=1
Sauer, Timothy (2017). \textit{Numerical Analysis}. 3rd ed. Pearson. \\
\hspace*{1.5em}\textsc{isbn}: 9780134696454

\noindent\hangindent=1.5em \hangafter=1
Krause, Jaden (2026). \textit{Initial Value Problems and Euler's Method}

\end{document}